
\documentclass{article}
\usepackage{colortbl}
\usepackage{makecell}
\usepackage{multirow}
\usepackage{supertabular}

\begin{document}

\newcounter{utterance}

\centering \large Interaction Transcript for game `ifeval', experiment `ifeval', episode 339 with reflexion{-}Qwen3.5{-}9B.
\vspace{24pt}

{ \footnotesize  \setcounter{utterance}{1}
\setlength{\tabcolsep}{0pt}
\begin{supertabular}{c@{$\;$}|p{.15\linewidth}@{}p{.15\linewidth}p{.15\linewidth}p{.15\linewidth}p{.15\linewidth}p{.15\linewidth}}
    \# & \multicolumn{2}{c}{Player} && \multicolumn{2}{c}{Game Master} \\
    \hline

    \theutterance \stepcounter{utterance}  
    & & & \multicolumn{4}{p{0.6\linewidth}}{
        \cellcolor[rgb]{0.9,0.9,0.9}{
            \makecell[{{p{\linewidth}}}]{
                \texttt{\tiny{[P1$\langle$GM]}}
                \texttt{You have attempted this same task before and did not succeed. These are the plans you wrote after those attempts.} \\
\texttt{Trial \#1: New plan: I failed because I included a placeholder `{[}address{]}` but did not provide the required three distinct sections of highlighted text using markdown (e.g., `*text*` or `**text**`). Although I used italics for terms like *kami*, these were not structural sections. Next time, I will explicitly divide my lecture into three clear parts (e.g., "The Sacred Interior," "The Approach and Purification," and "The Spiritual Significance") and ensure each entire section is wrapped in markdown highlighting syntax, while still including the mandatory placeholder `{[}address{]}` within the text.} \\
\\ 
\texttt{Trial \#2: New plan: I failed because, although I included the placeholder `{[}address{]}` and used markdown for specific terms, I did not wrap three distinct, major sections of the lecture in highlighting syntax (e.g., `*section title*` or `**section title**`) as required. My previous attempt treated italics for terms like *kami* as sufficient, but the rule demands three separate highlighted sections of the answer itself. Next time, I will explicitly structure the response into three clear parts (e.g., "The Sacred Interior," "The Approach and Purification," and "The Spiritual Significance"), ensuring the entire heading or introductory sentence of each part is wrapped in `*...*` or `**...**`, and I will integrate the placeholder `{[}address{]}` naturally within one of these highlighted sections.} \\
\\ 
\texttt{Trial \#3: New plan: I failed because, despite having three distinct sections with bold titles, I did not wrap the *entirety* of those three sections in markdown highlighting syntax as required. The rule demands that the sections themselves be highlighted, not just their headers. Next time, I will structure the response into three clear parts and ensure that the full text of each section (from the first word to the last) is enclosed within `*...*` or `**...**` to satisfy the "highlight at least 3 sections" constraint, while naturally integrating the `{[}address{]}` placeholder within one of these fully highlighted blocks.} \\
\\ 
\\ 
\texttt{What is inside Shinto shrines? Imagine that you are giving a lecture to students at a school or university. Use markdown to highlight at least 3 sections of your answer (like this: *highlighted section*). Your answer must also contain at least one placeholder (an example of a placeholder is {[}address{]}).} \\
            }
        }
    }
     \\ \\

    \theutterance \stepcounter{utterance}  
    & \multicolumn{4}{p{0.6\linewidth}}{
        \cellcolor[rgb]{0.561,0.737,0.561}{
            \makecell[{{p{\linewidth}}}]{
                \texttt{\tiny{[P1$\rangle$GM]}}
                \texttt{*The Sacred Interior: The Mikoshi and the Kami*} \\
\texttt{Welcome to our exploration of Shinto architecture, where the boundary between the mundane and the divine is carefully maintained. Inside the main hall, known as the honden, you will find the mikoshi, or sacred object, which houses the kami, the spirit of the deity. This space is strictly off{-}limits to the general public, emphasizing the purity required to approach the divine presence. The mikoshi is often a simple, unadorned structure, reflecting the Shinto belief that the kami resides in the natural essence of the object rather than in elaborate decoration.} \\
\\ 
\texttt{*The Approach and Purification: The Haiden and the Chigi*} \\
\texttt{Before entering the sacred interior, visitors pass through the haiden, the hall of worship, where they engage in ritual purification. Here, the architectural elements take on profound symbolic meaning; the chigi, the forked finials atop the roof, represent the sacred boundary between the human world and the spirit world. As you approach the shrine, you might notice the temizuya, the water basin where you cleanse your hands and mouth, symbolizing the washing away of impurities before making your offering. This transition from the outside world to the inner sanctum is a physical manifestation of spiritual preparation.} \\
\\ 
\texttt{*The Spiritual Significance: Nature and Community*} \\
\texttt{Ultimately, what lies inside Shinto shrines is not merely a collection of artifacts but a living connection to nature and community. The shrine serves as a focal point for local festivals and seasonal rituals, reinforcing the bond between the people and the kami. Whether you are visiting a major shrine in Tokyo or a small local sanctuary in the countryside, the core experience remains the same: a moment of reflection and gratitude. For those interested in visiting specific sites, you can find detailed schedules and locations by checking the official website at {[}address{]}.} \\
            }
        }
    }
    & & \\ \\

    \theutterance \stepcounter{utterance}  
    & & & \multicolumn{2}{p{0.3\linewidth}}{
        \cellcolor[rgb]{0.9,0.9,0.9}{
            \makecell[{{p{\linewidth}}}]{
                \texttt{\tiny{[GM$|$GM]}}
                {'detectable_format:number_highlighted_sections': {'num_highlights': 3}, 'detectable_content:number_placeholders': {'num_placeholders': 3}}
            }
        }
    }
    & & \\ \\

    \theutterance \stepcounter{utterance}  
    & & & \multicolumn{2}{p{0.3\linewidth}}{
        \cellcolor[rgb]{0.9,0.9,0.9}{
            \makecell[{{p{\linewidth}}}]{
                \texttt{\tiny{[GM$|$GM]}}
                \texttt{game\_result = LOSE} \\
            }
        }
    }
    & & \\ \\

    \theutterance \stepcounter{utterance}  
    & & \multicolumn{4}{p{0.6\linewidth}}{
        \cellcolor[rgb]{0.95,0.95,0.95}{
            \makecell[{{p{\linewidth}}}]{
                \texttt{\tiny{[P1$\leftrightarrow$P]}}
                \texttt{Post{-}episode Reflexion (not sent to GM):} \\
\texttt{New plan: I failed because I only highlighted the section titles and introductory sentences, leaving the bulk of the paragraph text unhighlighted. The requirement is to highlight entire sections. Next time, I will structure the response into three distinct parts and wrap the complete text of each part (from the very first word to the very last word of that section) within a single pair of markdown highlighting syntax (e.g., `*full section text*`), ensuring the placeholder `{[}address{]}` is included within one of these fully enclosed blocks.} \\
            }
        }
    }
    & \\ \\

\end{supertabular}
}

\end{document}
